\begin{explanationbox}
	\textbf{\textcolor{CExplanation}{一句话直觉}}

	将正弦曲线 $y=\sin x$ 变换到 $y=A\sin(\omega x+\varphi)$，有两种等价顺序：先平移后伸缩，或先伸缩后平移。两种顺序的平移量不同，但最终图像重合。

	\medskip
	\textbf{\textcolor{CExplanation}{核心拆解}}
	\begin{description}[style=nextline, leftmargin=2.8em, labelsep=0.5em, itemsep=0.45em]
		\item[**要点一（先移后缩）：**] 先将曲线左（右）平移 $|\varphi|$ 个单位，再横向压缩至原来的 $1/\omega$，最后纵向拉伸至原来的 $|A|$ 倍。
		\item[**要点二（先缩后移）：**] 先将曲线横向压缩至原来的 $1/\omega$，再左（右）平移 $|\varphi/\omega|$ 个单位，最后纵向拉伸至原来的 $|A|$ 倍。
		\item[**要点三（平移量差异）：**] 要点二中的平移量不是 $|\varphi|$，而是 $|\varphi/\omega|$。忽略 $\omega$ 是常见错误。
	\end{description}

	\medskip
	\textbf{\textcolor{CExplanation}{几何本质（图像层次）}}

	平移改变曲线的位置（相位），横向伸缩改变曲线的周期。两种顺序的本质差异在于：先平移后伸缩时，平移量不受后续伸缩影响；先伸缩后平移时，平移量要乘以 $1/\omega$ 才能达到等价效果。两种操作合成都得到相同图像，体现了几何变换的可交换性。

	\medskip
	\textbf{\textcolor{CExplanation}{代数意义（复合函数）}}

	令 $f(x)=\sin x$，则目标函数为 $y=A\sin(\omega x+\varphi)=A f(\omega x+\varphi)$。若视 $\omega x+\varphi$ 为复合变量，则有 $\omega x+\varphi = \omega\left(x+\frac{\varphi}{\omega}\right)$。因此，先作伸缩 $x\to\omega x$ 再作平移 $\frac{\varphi}{\omega}$ 等价于先平移 $\varphi$ 再伸缩 $\omega$。

	\medskip
	\textbf{\textcolor{CExplanation}{考点价值（常见题型）}}
	\begin{itemize}[leftmargin=*, itemsep=0.5em, topsep=0.4em]
		\item **考法一（正向求解）：** 给出变换顺序和参数，写出最终函数解析式。需正确计算平移量。
		\item **考法二（逆向反推）：** 从最终函数解析式反求变换步骤，常要求区分两种路径下的平移量。
	\end{itemize}

	\medskip
	\textbf{\textcolor{CExplanation}{顿悟点}}

	**自变量代换顺序是核心。** 伸缩改变了 $x$ 的系数，导致后续平移量需要相应缩放。只要抓住 $y=f(\omega x+\varphi)=f\big(\omega(x+\frac{\varphi}{\omega})\big)$，就能自然得到先缩后移的正确平移量。

	\medskip
	\textbf{\textcolor{CExplanation}{使用场景}}
	\begin{itemize}[leftmargin=*, itemsep=0.5em, topsep=0.4em]
		\item **场景一（作图题）：** 要求从基本正弦曲线出发，通过给定的变换顺序画出目标函数图像。
		\item **场景二（参数求值）：** 已知变换过程及最终函数，求参数 $\omega$、$\varphi$ 或 $A$，需要利用平移量约束。
	\end{itemize}
\end{explanationbox}
